\chapter{O sistema Sol--Terra--Lua}\label{ch:g6:sun-earth-moon}

Três mundos comandam o nosso céu: o Sol, que dá o dia; a Terra, que
nos carrega; a \omterm{def:g2:moon-in-the-sky:moon}{Lua}, que nos faz companhia. Você conhece cada um
deles; agora vamos mapear o \emph{sistema} --- quem gira em torno de
quem, com que rapidez, a que distância --- e pagar uma dívida antiga:
o segredo prometido das \omterm{def:g3:sun-days-seasons:year}{estações} está neste capítulo.

\section{Três movimentos, três relógios}

\begin{definition}[Eixo e órbita]\label{def:g6:sun-earth-moon:axisorbit}
O \emph{eixo}\index{eixo} da Terra é a linha imaginária que passa
pelos polos dela e em torno da qual ela gira, como o pino de um
pião. Uma \emph{órbita}\index{órbita} é o caminho fechado que um
corpo celeste percorre em torno de outro --- quase exatamente um
círculo para a Terra em torno do Sol e para a \omterm{def:g2:moon-in-the-sky:moon}{Lua} em torno da Terra.
\end{definition}

\begin{proposition}[Os três movimentos]\label{prop:g6:sun-earth-moon:motions}
O calendário inteiro se pendura em três movimentos constantes:
\begin{enumerate}
  \item a Terra \emph{gira} em torno do \omterm{def:g6:sun-earth-moon:axisorbit}{eixo} dela: uma volta em $24$
        \omterm{ex:g2:measuring-time:clock}{horas} --- o dia, desfilando o Sol e as \omterm{def:g4:solar-system:star}{estrelas} pelo nosso
        céu;
  \item a Terra \emph{orbita} o Sol: uma volta em cerca de $365$ dias
        e um quarto --- o ano, deslizando o arco do Sol para cima e
        para baixo ao longo das \omterm{def:g3:sun-days-seasons:year}{estações};
  \item a \omterm{def:g2:moon-in-the-sky:moon}{Lua} \emph{orbita} a Terra: uma volta em cerca de um mês ---
        e, como a metade iluminada dela fica voltada para \omterm{def:g5:series-parallel-circuits:parallel}{nós} aos
        poucos, desfilando as fases do seu antigo diário da \omterm{def:g2:moon-in-the-sky:moon}{Lua}.
\end{enumerate}
Três movimentos, três relógios: dia, mês, ano --- os medidores de
tempo mais antigos da humanidade, todos ainda funcionando.
\end{proposition}

\begin{omfigure}
\begin{tikzpicture}[scale=0.85]
  \fill[omThm] (5.6,2.2) circle (0.55);
  \foreach \a in {0,30,...,330} {
    \draw[omThm] (5.6,2.2) ++(\a:0.72) -- ++(\a:0.24);
  }
  \node[below, font=\small] at (5.6,1.4) {Sol};
  \draw[dashed, omExo, thick] (5.6,2.2) ellipse (4.6 and 1.75);
  \fill[omDef] (10.2,2.2) circle (0.24);
  \draw[-Stealth, thick] (10.2,2.75) arc (90:210:0.5);
  \node[above, font=\small] at (9.7,3.2) {giro: $24$ h};
  \draw[dashed, omExo] (10.2,2.2) circle (0.85);
  \fill[omExo] (10.8,2.8) circle (0.1);
  \node[right, font=\small] at (10.98,2.8) {Lua: um mês};
  \draw[-Stealth, thick, omDef] (8.1,0.62) arc (222:262:3.9);
  \node[below, font=\small, omDef] at (8.3,0.35) {órbita: um ano};
\end{tikzpicture}

{\small A planta do sistema (com as distâncias espremidas, como
sempre): a Terra gira enquanto orbita o Sol; a \omterm{def:g2:moon-in-the-sky:moon}{Lua} orbita a Terra.}
\end{omfigure}

\begin{example}[O quarto de dia e o ano bissexto]\label{ex:g6:sun-earth-moon:leap}
A volta da Terra leva $365$ dias \emph{e mais ou menos um quarto} ---
a natureza não deve lealdade nenhuma a números redondos. Os
calendários pagam a dívida guardando os quartos: quatro anos de
quartos de dia fazem um dia inteiro, acrescentado ao calendário a
cada quatro anos --- o \emph{ano bissexto}, de $366$ dias, com o dia
extra no fim de fevereiro. Deixe essa contabilidade de lado por um
século e o calendário se afasta semanas das \omterm{def:g3:sun-days-seasons:year}{estações} --- como
calendários mais antigos descobriram, dolorosamente.
\end{example}

\begin{example}[O sistema, em escala]\label{ex:g6:sun-earth-moon:scale}
Números para o caderno: a Terra tem cerca de \qty{12800}{km} de
diâmetro; a \omterm{def:g2:moon-in-the-sky:moon}{Lua}, cerca de um quarto disso, orbita a cerca de
\qty{384000}{km} de distância --- trinta Terras enfileiradas. O Sol
tem cerca de $109$ Terras de diâmetro e fica a $150$ milhões de
quilômetros. Encolha a Terra até virar uma bolinha de gude de um
\omterm{def:g2:measuring-length:units}{centímetro} e o modelo diz: a \omterm{def:g2:moon-in-the-sky:moon}{Lua}, um grão de pimenta a $30$
\omterm{def:g2:measuring-length:units}{centímetros}; o Sol, uma bola de mais de um \omterm{def:g2:measuring-length:units}{metro} de largura, parada a
mais de cem \omterm{def:g2:measuring-length:units}{metros} rua abaixo. A nossa vizinhança inteira é quase
toda vazio --- como o modelo em escala do \omterm{def:g4:solar-system:family}{Sistema Solar} na praça já
sussurrava.
\end{example}

\section{A dívida paga: por que existem estações}

\begin{proposition}[A inclinação faz as estações]\label{prop:g6:sun-earth-moon:tilt}
O \omterm{def:g6:sun-earth-moon:axisorbit}{eixo} da Terra não fica em pé sobre a \omterm{def:g6:sun-earth-moon:axisorbit}{órbita}: ele pende, por cerca
de um quarto de ângulo reto --- e continua pendendo \emph{para o
mesmo lado} durante toda a volta do ano, como uma \omterm{def:g4:magnets-and-compass:compass}{bússola} fiel das
\omterm{def:g4:solar-system:star}{estrelas}. Então, durante metade da volta, a nossa metade do mundo
pende \emph{na direção} do Sol: o Sol anda alto, os dias ficam
compridos --- verão. Meia volta depois, a mesma inclinação, sem
mudança nenhuma, nos aponta \emph{para longe} dele: Sol baixo, dias
curtos --- inverno. As \omterm{def:g3:sun-days-seasons:year}{estações} não são o Sol mudando, nem a
distância: são um \omterm{def:g4:solar-system:planet}{planeta} inclinado carregando a sua inclinação
imutável em volta da \omterm{def:g4:solar-system:star}{estrela} dele.
\end{proposition}

\begin{omfigure}
\begin{tikzpicture}[scale=0.85]
  \fill[omThm] (5.6,2.0) circle (0.5);
  \foreach \a in {0,45,...,315} {
    \draw[omThm] (5.6,2.0) ++(\a:0.65) -- ++(\a:0.22);
  }
  \begin{scope}[shift={(1.1,2.0)}]
    \draw[thick] (0,0) circle (0.62);
    \draw[very thick, omDef] (-0.28,-0.75) -- (0.28,0.75);
    \fill[omDef] (0.19,0.5) circle (0.09);
  \end{scope}
  \node[below, font=\small, align=center] at (1.1,0.9)
    {pendendo para o Sol:\\verão aqui};
  \begin{scope}[shift={(10.1,2.0)}]
    \draw[thick] (0,0) circle (0.62);
    \draw[very thick, omDef] (-0.28,-0.75) -- (0.28,0.75);
    \fill[omDef] (0.19,0.5) circle (0.09);
  \end{scope}
  \node[below, font=\small, align=center] at (10.1,0.9)
    {meia volta depois, mesma inclinação:\\pendendo para longe ---
     inverno aqui};
\end{tikzpicture}

{\small Uma mesma Terra inclinada em duas estações da volta dela. O
\omterm{def:g6:sun-earth-moon:axisorbit}{eixo} (azul) pende para o mesmo lado nos dois casos --- na direção do
Sol em junho, para longe dele em dezembro, para a cidade marcada no
hemisfério norte.}
\end{omfigure}

\begin{method}[O globo e a lâmpada]\label{met:g6:sun-earth-moon:globe}
Veja a máquina inteira sobre uma mesa:
\begin{enumerate}
  \item ponha uma lâmpada no meio da mesa como Sol e leve um globo
        inclinado (ou uma bola num espeto inclinado) devagar em volta
        dela, \emph{mantendo o espeto apontado para o mesmo canto da
        sala} o caminho inteiro;
  \item a cada quarto de volta, pare e olhe a latitude da sua cidade:
        veja quanto do círculo diário de giro dela fica na \omterm{def:g1:light-and-shadows:source}{luz};
  \item de um lado da lâmpada, o seu hemisfério se banha inclinado
        para dentro --- arcos iluminados compridos, lâmpada alta:
        verão; do lado oposto, inclinado para fora --- arcos curtos,
        lâmpada baixa: inverno;
  \item nas duas estações intermediárias, nenhum polo pende para
        dentro: o dia e a noite se dividem por igual --- primavera e
        outono.
\end{enumerate}
A única disciplina que faz isso funcionar: nunca deixe o espeto
girar --- a fidelidade do \omterm{def:g6:sun-earth-moon:axisorbit}{eixo} \emph{é} a razão das \omterm{def:g3:sun-days-seasons:year}{estações}.
\end{method}

\begin{example}[As velhas pistas, prestando contas]\label{ex:g6:sun-earth-moon:clues}
Cada observação dos seus anos de \omterm{def:g3:sun-days-seasons:sundial}{relógio de sol} agora se apresenta à
inclinação. O arco alto do verão: o nosso hemisfério pende para
dentro, então o Sol fica mais alto ao meio-dia. O arco atarracado e
as \omterm{def:g1:light-and-shadows:shadow}{sombras} compridas do inverno: pendendo para fora, o Sol raspa
baixo. \omterm{def:g3:sun-days-seasons:year}{Estações} opostas dos dois lados do equador --- férias de praia
em dezembro no sul enquanto o norte tira neve com a pá: quando um
hemisfério pende para dentro, o outro precisa pender para fora. Até
os dias iguais da primavera e do outono caem prontos da caminhada com
o globo e a lâmpada. Uma inclinação de um quarto de ângulo reto
comanda o espetáculo inteiro.
\end{example}

\begin{remark}[Não é a distância --- e a prova]\label{rem:g6:sun-earth-moon:notdistance}
A \omterm{def:g6:sun-earth-moon:axisorbit}{órbita} da Terra é tão próxima de um círculo que a distância até o
Sol muda pouco --- e, para nosso deleite, no inverno do hemisfério
norte estamos ligeiramente \emph{mais perto} do Sol. As \omterm{def:g3:sun-days-seasons:year}{estações} se
importam com a \emph{inclinação} da \omterm{def:g1:light-and-shadows:source}{luz}, e não com o comprimento da
viagem: pendendo para dentro, a \omterm{def:g1:light-and-shadows:source}{luz} do Sol bate íngreme e
concentrada; pendendo para fora, a mesma \omterm{def:g1:light-and-shadows:source}{luz} se espalha fina pelo
chão. Uma lanterna apontada de frente e depois inclinada contra uma
parede mostra a diferença num segundo --- a \omterm{def:g1:light-and-shadows:source}{luz} íngreme morde, a \omterm{def:g1:light-and-shadows:source}{luz}
rasante mal aquece.
\end{remark}

\section{Exercícios}

\begin{exercise}[$\star$]\label{exo:g6:sun-earth-moon:1}
Diga os três movimentos do sistema e o relógio que cada um deles dá
corda.
\end{exercise}

\begin{exercise}[$\star$]\label{exo:g6:sun-earth-moon:2}
O que é o \omterm{def:g6:sun-earth-moon:axisorbit}{eixo} da Terra? Qual fato de um quarto de ângulo reto sobre
ele, mais qual fidelidade, produz as \omterm{def:g3:sun-days-seasons:year}{estações}?
\end{exercise}

\begin{exercise}[$\star$]\label{exo:g6:sun-earth-moon:3}
Por que o calendário acrescenta um dia a cada quatro anos? O que
escorregaria se ele nunca fizesse isso?
\end{exercise}

\begin{exercise}[$\star$]\label{exo:g6:sun-earth-moon:4}
Aproximadamente quantas Terras enfileiradas chegam à \omterm{def:g2:moon-in-the-sky:moon}{Lua}? Quantas
Terras cobrem a face do Sol?
\end{exercise}

\begin{exercise}[$\star$]\label{exo:g6:sun-earth-moon:5}
Na caminhada com o globo e a lâmpada, por que o espeto precisa
continuar apontando para o mesmo canto da sala? O que se quebra se
ele girar?
\end{exercise}

\begin{exercise}[$\star$]\label{exo:g6:sun-earth-moon:6}
Quando é verão aqui, que \omterm{def:g3:sun-days-seasons:year}{estação} é no hemisfério do outro lado do
equador, e por que tem de ser assim?
\end{exercise}

\begin{exercise}[$\star$]\label{exo:g6:sun-earth-moon:7}
Use a \omterm{prop:g5:mirrors-reflection:image}{imagem} da lanterna contra a parede da
\cref{rem:g6:sun-earth-moon:notdistance} para explicar por que
pender na direção do Sol nos aquece mais.
\end{exercise}

\begin{exercise}[$\star$]\label{exo:g6:sun-earth-moon:8}
Uma velha conhecida volta: por que a \omterm{def:g2:moon-in-the-sky:moon}{Lua} mostra fases ao longo do
mês? Qual dos três movimentos comanda esse espetáculo?
\end{exercise}

\begin{exercise}[$\star\star$]\label{exo:g6:sun-earth-moon:9}
Um tio teimoso insiste: ``O inverno é quando a Terra está mais longe
do Sol.'' Reúna dois fatos deste capítulo que, juntos, o derrotam.
\end{exercise}

\begin{exercise}[$\star\star$]\label{exo:g6:sun-earth-moon:10}
Suponha que o \omterm{def:g6:sun-earth-moon:axisorbit}{eixo} da Terra ficasse perfeitamente em pé sobre a
\omterm{def:g6:sun-earth-moon:axisorbit}{órbita} --- sem inclinação nenhuma. Descreva um ano: o que acontece
com as \omterm{def:g3:sun-days-seasons:year}{estações}, com a altura do Sol ao meio-dia e com a \omterm{def:g2:measuring-time:duration}{duração} dos
dias? (Faça o globo em pé dar a volta na lâmpada dentro da sua
cabeça.)
\end{exercise}

\begin{exercise}[$\star\star$]\label{exo:g6:sun-earth-moon:11}
A volta da \omterm{def:g2:moon-in-the-sky:moon}{Lua} leva cerca de $27$ dias e, mesmo assim, de \omterm{ex:g2:moon-in-the-sky:shapes}{lua cheia} a
\omterm{ex:g2:moon-in-the-sky:shapes}{lua cheia} passam cerca de $29.5$. Proponha a solução --- o que mais
se moveu enquanto a \omterm{def:g2:moon-in-the-sky:moon}{Lua} dava a volta? (Esboce a Terra avançando pela
\omterm{def:g6:sun-earth-moon:axisorbit}{órbita} dela durante uma volta da \omterm{def:g2:moon-in-the-sky:moon}{Lua}.)
\end{exercise}

\begin{exercise}[$\star\star\star$]\label{exo:g6:sun-earth-moon:12}
No auge do verão, as regiões perto do Polo Norte mantêm o Sol acima
do horizonte à meia-noite --- o famoso ``sol da meia-noite'' --- e no
auge do inverno elas esperam semanas por um \omterm{def:g1:day-and-night:daynight}{nascer do sol}. Explique
as duas maravilhas com os círculos de giro inclinados do
\cref{met:g6:sun-earth-moon:globe}: o que a inclinação faz com o
círculo diário de uma cidade polar em junho, e em dezembro?
\end{exercise}

\section{Problema: O modelo no campo de esportes}

\begin{problem}[{Problema de fim de semana --- a turma constrói o
sistema Sol--Terra--Lua no campo de esportes; bolinhas de gude, grãos
de pimenta e uma caminhada bem longa; a inclinação entra no
modelo}]\label{pb:g6:sun-earth-moon:1}
A turma constrói o sistema em escala: a Terra é uma bolinha de gude
de \qty{1}{cm}. Números verdadeiros úteis: Terra com \qty{12800}{km}
de diâmetro; \omterm{def:g2:moon-in-the-sky:moon}{Lua} com \qty{3200}{km} de diâmetro, orbitando a
\qty{384000}{km}; Sol com $109$ Terras de diâmetro, a $150$ milhões
de quilômetros.

\medskip\noindent\textbf{Parte I --- Escalando os papéis.}
\begin{enumerate}
  \item No modelo, \qty{1}{cm} representa quantos quilômetros?
  \item A \omterm{def:g2:moon-in-the-sky:moon}{Lua} tem cerca de um quarto da largura da Terra. Qual é a
        largura da \omterm{def:g2:moon-in-the-sky:moon}{Lua} do modelo --- e qual item de cozinha faz bem
        esse papel?
  \item A distância verdadeira da \omterm{def:g2:moon-in-the-sky:moon}{Lua} é de $30$ larguras da Terra.
        Onde fica a \omterm{def:g2:moon-in-the-sky:moon}{Lua} do modelo, a partir da bolinha de gude?
  \item O Sol do modelo tem $109$ bolinhas de largura. Quanto é isso
        em \omterm{def:g2:measuring-length:units}{centímetros} --- e qual objeto de pátio tem mais ou menos
        esse tamanho?
\end{enumerate}

\medskip\noindent\textbf{Parte II --- A caminhada longa.}
\begin{enumerate}[resume]
  \item A distância do Sol é de $150$ milhões de km. No modelo ---
        divida pela resposta da questão 1 --- quanto é isso em
        \omterm{def:g2:measuring-length:units}{centímetros}, e em \omterm{def:g2:measuring-length:units}{metros}?
  \item O modelo cabe num campo de esportes de \qty{100}{m}? Onde a
        bola-Sol precisa ficar?
  \item Um aluno percorre o trecho Terra--Sol a $1$ \omterm{def:g2:measuring-length:units}{metro} por
        segundo. A caminhada leva mais ou menos quantos minutos?
  \item Parada na bolinha de gude e olhando para trás, para a
        bola-Sol de um \omterm{def:g2:measuring-length:units}{metro} a $117$ \omterm{def:g2:measuring-length:units}{metros} de distância, a turma se
        surpreende: ela parece ter mais ou menos o mesmo tamanho do
        \omterm{def:g2:moon-in-the-sky:moon}{grão-Lua} de um quarto de \omterm{def:g2:measuring-length:units}{centímetro} que está a $30$
        \omterm{def:g2:measuring-length:units}{centímetros} da bolinha. Confira as duas razões --- e ligue
        essa coincidência a algo que acontece, raramente e de forma
        gloriosa, no céu de verdade.
\end{enumerate}

\medskip\noindent\textbf{Parte III --- Fazendo o modelo funcionar.}
\begin{enumerate}[resume]
  \item Para animar o modelo com honestidade, quanto tempo a bolinha
        deve levar para dar a volta na bola-Sol, e o grão para dar a
        volta na bolinha, se um dia real for representado por um
        segundo? (Conte por alto: o ano como $365$ dias e a volta da
        \omterm{def:g2:moon-in-the-sky:moon}{Lua} como $27$ dias.)
  \item A bolinha também precisa girar em torno de si mesma. A um dia
        real por segundo, com que rapidez isso é no modelo --- e por
        que ninguém vai enxergar esse giro numa bolinha de
        \qty{1}{cm}?
  \item Um espeto atravessando a bolinha faz o papel do \omterm{def:g6:sun-earth-moon:axisorbit}{eixo}. Com que
        inclinação, e apontando para onde, o espeto precisa ser
        carregado ao longo da volta inteira?
  \item Escreva o aviso honesto do modelo: diga uma coisa que o
        modelo do campo de esportes mostra com verdade e uma coisa
        que nenhum modelo do tamanho de um campo consegue mostrar ao
        mesmo tempo (tamanhos e distâncias juntos? os movimentos? a
        \omterm{def:g1:light-and-shadows:source}{luz}?). Justifique com os seus números.
\end{enumerate}
\end{problem}
